%! Introducing Vectors — Unit 1, Lesson 1.0 — Group Activity (Answer Key)
\documentclass[10pt]{article}
\usepackage{linalg-article}
\usepackage{linalg-key}   % pulls in -boxes; do NOT also load -boxes
\newcommand{\TallMath}[1]{$\displaystyle #1\rule[-1.4em]{0pt}{3.2em}$}
\newcommand{\vv}{\mathbf{v}}
\newcommand{\ww}{\mathbf{w}}

\begin{document}

\pageheader{Unit 1, Lesson 1.0}{Group Activity --- Answer Key}
\namepartnerperiod

\vspace{0.10in}

\begin{headlinebox}{blush}
\small\textbf{Drone Delivery.} A delivery drone moves over a city grid. Each hop is a
vector $\begin{bsmallmatrix} \text{east} \\ \text{north} \end{bsmallmatrix}$ measured in
blocks (west and south are negative). Work with your partner and \emph{show your steps}.
\end{headlinebox}

\vspace{0.10in}

\begin{tcolorbox}[colback=white, colframe=black!40, title=\textbf{Tier R --- Remediate},
    fonttitle=\bfseries, arc=1.5mm, left=3mm, right=3mm, top=2mm, bottom=2mm, breakable]
The drone's first two hops are $\vv = \begin{bmatrix} 4 \\ 3 \end{bmatrix}$ and
$\ww = \begin{bmatrix} -1 \\ 2 \end{bmatrix}$.
\begin{enumerate}[leftmargin=1.7em, itemsep=8pt]
  \item In words, what does hop $\ww$ tell the drone to do? \ansline{Move $1$ block west and $2$ blocks north}
  \item Find $\vv + \ww$ (the two hops combined). \ans{$\begin{bmatrix} 3 \\ 5 \end{bmatrix}$}
  \item Find $2\vv$. What happens to the arrow when you double it?
        \ans{$\begin{bmatrix} 8 \\ 6 \end{bmatrix}$; twice as long, same direction.}
  \item Find $\|\vv\|$, the length of the first hop. Show the Pythagorean step.
        \ansline{$\sqrt{4^2+3^2}=\sqrt{25}=5$}
\end{enumerate}
\end{tcolorbox}

\begin{tcolorbox}[colback=white, colframe=black!40, title=\textbf{Tier A --- Approaching Proficiency},
    fonttitle=\bfseries, arc=1.5mm, left=3mm, right=3mm, top=2mm, bottom=2mm, breakable]
A full delivery is three hops from the depot at the origin:
\[
  \mathbf{a} = \begin{bmatrix} 6 \\ 0 \end{bmatrix}, \quad
  \mathbf{b} = \begin{bmatrix} 2 \\ 6 \end{bmatrix}, \quad
  \mathbf{c} = \begin{bmatrix} -2 \\ 2 \end{bmatrix}.
\]
\begin{enumerate}[leftmargin=1.7em, itemsep=8pt]
  \item Find the drone's final position $\mathbf{a}+\mathbf{b}+\mathbf{c}$.
        \ans{$\begin{bmatrix} 6 \\ 8 \end{bmatrix}$}
  \item How far east and how far north of the depot does it end up?
        \ansline{$6$ blocks east and $8$ blocks north}
  \item Find the straight-line distance from the depot to the final position.
        \ansline{$\sqrt{6^2+8^2}=\sqrt{100}=10$ blocks}
  \item The drone flew the three hops but the \emph{customer} only cares about the
        straight-line distance. Explain why those two totals are different.
        \ansline{Path flown $\approx 6+6.3+2.8 = 15.1$ blocks, but the straight-line
        distance is only $10$. The straight line is the shortest route; changing direction
        adds length.}
\end{enumerate}
\end{tcolorbox}

\begin{tcolorbox}[colback=white, colframe=black!40, title=\textbf{Tier E --- Extension},
    fonttitle=\bfseries, arc=1.5mm, left=3mm, right=3mm, top=2mm, bottom=2mm, breakable]
\begin{enumerate}[leftmargin=1.7em, itemsep=8pt]
  \item A hop $\vv=\begin{bmatrix} 8 \\ 6 \end{bmatrix}$ has length $10$. Find a
        \textbf{unit vector} pointing the same direction as $\vv$.
        \ansline{$\frac{1}{10}\begin{bsmallmatrix} 8 \\ 6 \end{bsmallmatrix}
        = \begin{bsmallmatrix} 0.8 \\ 0.6 \end{bsmallmatrix}$;\ \ check: $\sqrt{0.8^2+0.6^2}=1$}
  \item \textbf{Justify.} For two hops $\vv$ and $\ww$, is it always true that
        $\|\vv+\ww\| = \|\vv\| + \|\ww\|$? Test it on
        $\vv=\begin{bmatrix} 3 \\ 0\end{bmatrix},\ \ww=\begin{bmatrix} 0 \\ 4 \end{bmatrix}$,
        then explain when the two sides are equal.
        \ansline{$\|\vv+\ww\|=\|(3,4)\|=5$ but $\|\vv\|+\|\ww\|=3+4=7$, so \emph{not} always.}
        \ansline{They are equal only when $\vv$ and $\ww$ point the same direction (one is a
        positive multiple of the other). Otherwise $\|\vv+\ww\|<\|\vv\|+\|\ww\|$ --- the
        triangle inequality.}
\end{enumerate}
\end{tcolorbox}

\begin{teachernote}
Tier A item 4 is the key ``interpret'' moment: separate \emph{distance travelled} (sum of
lengths) from \emph{displacement} (length of the sum). Tier E item 2 previews the triangle
inequality; accept an informal ``a straight line beats a detour'' argument.
\end{teachernote}

\end{document}
